\section{ФСР однородной системы. Теорема о ФСР. Примеры нахождения ФСР}

\begin{definition}
\textbf{Фундаментальной системой решений} (ФСР) однородной системы линейных уравнений $Ax = 0$ называется любая система решений $y^{(1)}, y^{(2)}, \dots, y^{(k)}$, удовлетворяющая следующим условиям:
\begin{enumerate}
    \item Каждый вектор $y^{(i)}$ является решением системы $Ax = 0$.
    \item Система векторов $\{y^{(1)}, y^{(2)}, \dots, y^{(k)}\}$ линейно независима.
    \item Любое решение $X$ системы $Ax = 0$ линейно выражается через эту систему векторов.
\end{enumerate}
\end{definition}

\begin{theorem}[О структуре общего решения однородной системы]
Пусть дана однородная система $Ax = 0$ с $n$ неизвестными, и $\operatorname{rang} A = r$. 
\begin{enumerate}
    \item Если $r = n$, то система имеет только тривиальное решение $X = 0$.
    \item Если $r < n$, то существует ФСР, состоящая ровно из $n - r$ решений $y^{(1)}, \dots, y^{(n-r)}$, и общее решение системы имеет вид:
    \[
    X_{оо} = c_1 y^{(1)} + c_2 y^{(2)} + \dots + c_{n-r} y^{(n-r)},
    \]
    где $c_1, \dots, c_{n-r}$ --- произвольные постоянные.
\end{enumerate}
\end{theorem}

\begin{proof}
Приведем расширенную матрицу системы к ступенчатому виду. Поскольку система однородна, столбец свободных членов состоит из нулей и не влияет на преобразования. После приведения к ступенчатому виду система имеет $r$ ненулевых строк. 
Выделим $r$ \textit{главных} переменных (соответствующих ведущим элементам) и $n - r$ \textit{свободных} переменных. Перенесем слагаемые со свободными переменными в правые части:
\[
\begin{cases}
a_{11}x_1 + \dots + a_{1r}x_r = -a_{1,r+1}x_{r+1} - \dots - a_{1n}x_n \\
\vdots \\
a_{rr}x_r = -a_{r,r+1}x_{r+1} - \dots - a_{rn}x_n
\end{cases}
\]
Разделим каждое уравнение на его ведущий коэффициент, чтобы выразить главные переменные через свободные:
\[
\begin{cases}
x_1 = c_{1,r+1}x_{r+1} + \dots + c_{1n}x_n \\
x_2 = c_{2,r+1}x_{r+1} + \dots + c_{2n}x_n \\
\vdots \\
x_r = c_{r,r+1}x_{r+1} + \dots + c_{rn}x_n
\end{cases}
\]
Построим $n - r$ частных решений, придавая свободным переменным $(x_{r+1}, \dots, x_n)$ значения строк единичной матрицы порядка $n - r$:
\[
\begin{array}{c|cccc}
\text{Решение} & x_{r+1} & x_{r+2} & \dots & x_n \\
\hline
y^{(1)} & 1 & 0 & \dots & 0 \\
y^{(2)} & 0 & 1 & \dots & 0 \\
\vdots & \vdots & \vdots & \ddots & \vdots \\
y^{(n-r)} & 0 & 0 & \dots & 1
\end{array}
\]
Подставляя эти значения, мы однозначно находим соответствующие значения главных переменных. Полученные векторы $y^{(1)}, \dots, y^{(n-r)}$ являются решениями системы. Они линейно независимы, так как их последние $n - r$ координат образуют единичную матрицу (и любая их линейная комбинация, равная нулю, потребует равенства нулю всех коэффициентов).

Любое решение $X$ системы получается при некотором выборе свободных переменных $x_{r+1}=c_1, \dots, x_n=c_{n-r}$. В силу линейности выражений для главных переменных, этот вектор $X$ в точности равен $c_1 y^{(1)} + \dots + c_{n-r} y^{(n-r)}$. Таким образом, построенная система является ФСР.
\hfill$\blacksquare$
\end{proof}

\begin{example}[Пример нахождения ФСР]
Пусть после приведения к ступенчатому виду система имеет вид:
\[
\begin{cases}
x_1 + 2x_3 - x_4 = 0 \\
x_2 - x_3 + 3x_4 = 0
\end{cases}
\]
Здесь $n=4$, $r=2$. Главные переменные: $x_1, x_2$. Свободные переменные: $x_3, x_4$.
Выражаем главные через свободные:
\[
\begin{cases}
x_1 = -2x_3 + x_4 \\
x_2 = x_3 - 3x_4
\end{cases}
\]
Составляем матрицу решений, где нижний блок --- единичная матрица для свободных переменных, а верхний блок --- коэффициенты при них:
\[
\begin{pmatrix} x_1 \\ x_2 \\ x_3 \\ x_4 \end{pmatrix} = 
x_3 \begin{pmatrix} -2 \\ 1 \\ 1 \\ 0 \end{pmatrix} + 
x_4 \begin{pmatrix} 1 \\ -3 \\ 0 \\ 1 \end{pmatrix}
\]
Следовательно, ФСР состоит из векторов $y^{(1)} = (-2, 1, 1, 0)^T$ и $y^{(2)} = (1, -3, 0, 1)^T$.
\end{example}
